\chapter{Basis Data}
\label{ch:07-basis-data}

\section{Engine dan Keputusan Desain}

Basis data sistem ini adalah SQLite yang berjalan di dalam proses aplikasi, dengan
berkas tunggal \code{\DbFile}. Tidak ada server basis data terpisah dan tidak ada
layanan yang harus dinyalakan lebih dahulu. Pilihan itu diambil karena server
mandiri seperti MySQL memerlukan proses terpisah, sedangkan SQLite dapat dimuat
sebagai pustaka di dalam proses Python yang juga menjalankan broker, worker, dan
dashboard.

Skema pada \file{database/schema.sql} berisi \NumTables{} tabel dan satu view
agregat, dan seluruh definisinya idempoten sehingga dapat dijalankan pada basis data
kosong maupun yang sudah terisi. Tiga penerjemahan dari dialek server ke SQLite
dipakai konsisten: enum diganti kolom teks dengan \code{CHECK}, JSON diganti kolom
teks dengan \code{json_valid()}, dan waktu berpresisi milidetik diganti kolom teks
yang diisi \code{strftime}. Seluruh waktu disimpan dalam UTC dan dikonversi ke waktu
lokal di lapisan aplikasi.

Konkurensi ditangani dengan tiga pengaturan pada setiap koneksi: mode jurnal WAL
agar satu penulis dan banyak pembaca dapat bekerja bersamaan, \code{busy_timeout}
lima detik agar penulis menunggu alih-alih langsung gagal, dan \code{foreign_keys}
yang diaktifkan eksplisit karena SQLite mematikannya secara bawaan. Setiap thread
memakai koneksi sendiri, dan operasi yang menulis lebih dari satu baris dibungkus
\code{BEGIN IMMEDIATE}. (\file{dashboard/db.py})

\section{Diagram Relasi}

\figref{fig:db-er} memuat diagram relasi antar entitas, sedangkan
\figref{fig:db-relational-schema} menampilkan skema relasional dengan kunci primer,
kunci asing, dan index. Yang pertama dipakai untuk memahami kardinalitas, yang kedua
untuk memeriksa kolom.

\figureasset[\linewidth]{db-er}{Diagram relasi entitas: perangkat sebagai pusat, dengan
telemetri, jadwal, perintah, riwayat pakan, dan notifikasi sebagai entitas
turunan.}{fig:db-er}

\figuretikz[\linewidth]{db-relational-schema}{Skema relasional basis data: \NumTables{} tabel,
kunci primer, kunci asing, dan index.}{fig:db-relational-schema}

\section{Kamus Data}

\tabref{tab:07-kamus} memuat seluruh kolom pada keenam tabel. Kolom ``Constraint''
menuliskan batasan tingkat basis data, bukan aturan yang ditegakkan kode aplikasi;
aturan yang hanya ada di kode dibahas pada bagian terakhir bab ini.

{\scriptsize
\begin{longtable}{@{}p{2.1cm}p{3.8cm}p{1.5cm}p{2.5cm}p{4.1cm}@{}}
\caption{Kamus data basis data.}
\label{tab:07-kamus} \\
\toprule
\textbf{Tabel} & \textbf{Kolom} & \textbf{Tipe} & \textbf{Constraint} &
\textbf{Keterangan} \\
\midrule
\endfirsthead
\toprule
\textbf{Tabel} & \textbf{Kolom} & \textbf{Tipe} & \textbf{Constraint} &
\textbf{Keterangan} \\
\midrule
\endhead
\bottomrule
\endfoot
devices & id & INTEGER & primary key & Kunci utama. \\
 & \code{device_code} & TEXT & not null, unik & Kode perangkat, segmen topik MQTT. \\
 & name & TEXT & not null & Nama tampilan. \\
 & location & TEXT & boleh null & Lokasi penempatan. \\
 & timezone & TEXT & not null, bawaan Asia/Jakarta & Zona untuk menyetel RTC. \\
 & \code{sensor_full_cm} & REAL & not null, bawaan \HopperFull & Jarak saat hopper
penuh. \\
 & \code{sensor_empty_cm} & REAL & not null, bawaan \HopperEmpty & Jarak saat hopper
kosong. \\
 & \code{low_feed_threshold_pct} & INTEGER & 0 sampai 100, bawaan \DefaultThreshold
& Ambang pakan menipis. \\
 & \code{is_online} & INTEGER & not null, 0 atau 1 & Dari topik status. \\
 & \code{last_seen_at} & TEXT & boleh null & Waktu data terakhir. \\
 & \code{created_at}, \code{updated_at} & TEXT & not null, waktu sekarang & Diisi
kode. \\
\midrule
telemetry & id & INTEGER & primary key & Kunci utama. \\
 & \code{device_id} & INTEGER & kunci asing, cascade & Menunjuk ke \code{devices}. \\
 & \code{distance_cm} & REAL & boleh null & Jarak mentah sensor. \\
 & \code{feed_level_pct} & REAL & boleh null & Persentase sisa pakan. \\
 & state & TEXT & boleh null & \code{idle}, \code{feeding}, \code{error}. \\
 & \code{recorded_at} & TEXT & not null & Waktu baris masuk. \\
\midrule
schedules & id & INTEGER & primary key & Kunci utama. \\
 & \code{device_id} & INTEGER & kunci asing, cascade & Menunjuk ke \code{devices}. \\
 & name & TEXT & not null & Nama jadwal. \\
 & \code{feed_time} & TEXT & not null & Jam pemberian, misalnya 07:00:00. \\
 & portions & INTEGER & not null, bawaan 1 & Porsi per eksekusi. \\
 & \code{is_active} & INTEGER & not null, 0 atau 1 & Hanya jadwal aktif yang dikirim. \\
 & \code{created_at}, \code{updated_at} & TEXT & not null & Waktu pembuatan dan
perubahan. \\
\midrule
commands & id & INTEGER & primary key & Kunci utama. \\
 & \code{command_id} & TEXT & not null, unik & Pengenal dari payload. \\
 & \code{device_id} & INTEGER & kunci asing, cascade & Menunjuk ke \code{devices}. \\
 & type & TEXT & not null, enam nilai & Enam tipe pada \chapref{ch:06-protokol-mqtt}. \\
 & payload & TEXT & harus JSON valid & Payload apa adanya. \\
 & status & TEXT & bawaan pending & \code{sent}, \code{success}, \code{failed},
\code{timeout}. \\
 & \code{result_message} & TEXT & boleh null & Pesan dari ack. \\
 & \code{created_at} & TEXT & not null & Waktu perintah dibuat. \\
 & \code{acked_at} & TEXT & boleh null & Waktu ack diterima. \\
\midrule
\code{feeding_logs} & id & INTEGER & primary key & Kunci utama. \\
 & \code{device_id} & INTEGER & kunci asing, cascade & Menunjuk ke \code{devices}. \\
 & \code{command_id} & TEXT & boleh null & Korelasi perintah; null untuk jadwal. \\
 & \code{trigger_type} & TEXT & \code{manual} atau \code{scheduled} & Pemicu
pemberian pakan. \\
 & \code{schedule_id} & INTEGER & kunci asing, set null & Selalu null saat runtime. \\
 & portions & INTEGER & not null, bawaan 1 & Jumlah porsi. \\
 & status & TEXT & bawaan pending & \code{pending}, \code{success}, \code{failed}. \\
 & message & TEXT & boleh null & Pesan hasil dari ack. \\
 & \code{distance_before_cm}, \code{distance_after_cm} & REAL & boleh null & Jarak
sebelum dan sesudah. \\
 & \code{feed_level_before_pct}, \code{feed_level_after_pct} & REAL & boleh null &
Persentase sebelum dan sesudah. \\
 & \code{requested_at} & TEXT & not null & Waktu permintaan. \\
 & \code{completed_at} & TEXT & boleh null & Waktu selesai. \\
\midrule
notifications & id & INTEGER & primary key & Kunci utama. \\
 & \code{device_id} & INTEGER & kunci asing, cascade & Menunjuk ke \code{devices}. \\
 & type & TEXT & not null, empat nilai & \code{feed_low}, \code{feed_empty},
\code{feed_failed}, \code{device_offline}. \\
 & severity & TEXT & bawaan warning & \code{info}, \code{warning},
\code{critical}. \\
 & title & TEXT & boleh null & Judul singkat. \\
 & message & TEXT & not null & Uraian kejadian. \\
 & \code{feed_level_pct} & REAL & boleh null & Level saat notifikasi dibuat. \\
 & \code{is_read} & INTEGER & not null, 0 atau 1 & Penanda sudah dibaca. \\
 & \code{created_at} & TEXT & not null & Waktu notifikasi dibuat. \\
\end{longtable}
}

\src{database/schema.sql}

\section{Index dan Batasan Unik}

\tabref{tab:07-index} memuat index eksplisit beserta alasannya. Selain itu ada dua
batasan unik yang otomatis dibuatkan index oleh SQLite, yaitu pada
\code{devices.device_code} dan \code{commands.command_id}. Keduanya bukan sekadar
validasi: \code{device_code} adalah identitas perangkat pada topik MQTT, dan
\code{command_id} adalah kunci korelasi ack sehingga duplikat harus ditolak pada
tingkat basis data.

\begin{table}[H]
\centering
\caption{Index eksplisit pada basis data.}
\label{tab:07-index}
\scriptsize
\begin{tabular}{@{}p{5.0cm}p{1.9cm}p{2.3cm}p{4.6cm}@{}}
\toprule
\textbf{Index} & \textbf{Tabel} & \textbf{Kolom} & \textbf{Alasan} \\
\midrule
\code{idx_telemetry_device_time} & telemetry & device, waktu & Grafik rentang waktu
dan baris terakhir. \\
\code{idx_schedules_device_active} & schedules & device, aktif, jam & Daftar jadwal
aktif terurut jam. \\
\code{idx_commands_device_time} & commands & device, waktu & Aktivitas terakhir dan
sapuan timeout. \\
\code{idx_feeding_device_time} & \code{feeding_logs} & device, waktu & Riwayat pakan
per periode. \\
\code{idx_feeding_status} & \code{feeding_logs} & status & Riwayat yang belum
selesai. \\
\code{idx_notifications_device} & notifications & device, dibaca, waktu & Notifikasi
belum dibaca. \\
\bottomrule
\end{tabular}
\end{table}

\section{View Agregat}

View \code{v_device_status} menyatukan satu baris perangkat dengan baris telemetri
terbarunya; baris terbaru dipilih dengan fungsi jendela \code{ROW_NUMBER} yang
diurutkan menurun menurut waktu catat dan nomor baris, sehingga dua baris bercap
waktu sama tidak menghasilkan duplikat. View ini juga menghitung dua penanda yang
dipakai dashboard: \code{is_empty} bernilai 1 bila persentase nol atau kurang dan
null bila persentase belum diketahui, sedangkan \code{is_low} bernilai 1 bila
persentase di atas nol tetapi tidak lebih dari ambang perangkat. Pembedaan ``belum
diketahui'' dan ``kosong'' dipertahankan dengan sengaja: perangkat yang sensornya
rusak mengirim telemetri tanpa persentase, dan view tidak boleh melaporkannya
sebagai hopper kosong. Agregat ini tidak disimpan sebagai tabel agar tidak ada
salinan yang dapat menyimpang dari data aslinya.

\section{Data Awal}

Basis data baru diisi dua hal: satu baris perangkat \code{\SeedDevice} bernama
\SeedDeviceName{} dengan kalibrasi \HopperFull{} cm dan \HopperEmpty{} cm serta
ambang \DefaultThreshold{} persen, dan dua jadwal contoh untuk perangkat itu, yaitu
pakan pagi pukul 07:00 dan pakan sore pukul 17:00, masing-masing dua porsi.
Penyisipan memakai pola yang aman dijalankan berulang: baris perangkat memakai
\code{INSERT OR IGNORE}, sedangkan jadwal hanya disisipkan bila belum ada jadwal
pada jam yang sama. Dengan pola itu, menjalankan skema pada basis data berisi data
pengguna tidak mengubah apa pun.

\section{Kardinalitas}

\tabref{tab:07-kardinalitas} merangkum hubungan antar tabel dan perilakunya saat
baris induk dihapus. Semua hubungan berasal dari perangkat dan memakai cascade,
sehingga menghapus satu perangkat menghapus seluruh jejaknya. Satu-satunya hubungan
tanpa cascade adalah riwayat pakan ke jadwal, yang memakai \code{SET NULL} agar
riwayat tidak hilang ketika jadwal dihapus.

\begin{table}[H]
\centering
\caption{Kardinalitas antar tabel dan perilaku saat induk dihapus.}
\label{tab:07-kardinalitas}
\small
\begin{tabular}{@{}p{4.3cm}p{2.2cm}p{2.3cm}p{5.6cm}@{}}
\toprule
\textbf{Hubungan} & \textbf{Kardinalitas} & \textbf{Saat induk dihapus} &
\textbf{Catatan} \\
\midrule
devices ke telemetry & satu ke banyak & cascade & Satu baris per sampel sensor. \\
devices ke schedules & satu ke banyak & cascade & Jadwal berlaku setiap hari. \\
devices ke commands & satu ke banyak & cascade & Termasuk perintah yang gagal. \\
devices ke \code{feeding_logs} & satu ke banyak & cascade & Riwayat manual dan
terjadwal. \\
devices ke notifications & satu ke banyak & cascade & Notifikasi tidak berdiri
sendiri. \\
schedules ke \code{feeding_logs} & satu ke nol atau banyak & set null & Selalu null
saat runtime. \\
commands ke \code{feeding_logs} & satu ke nol atau satu & tanpa kunci asing &
Korelasi lewat \code{command_id}. \\
devices ke view status & satu ke satu & tidak berlaku & View, bukan tabel. \\
\bottomrule
\end{tabular}
\end{table}

\section{Migrasi Saat Inisialisasi}

Inisialisasi tidak berhenti pada pembuatan skema. Ada tiga langkah pemeliharaan yang
dijalankan setiap kali proses dinyalakan dan semuanya idempoten: kolom
\code{timezone} ditambahkan bila belum ada; baris perangkat yang masih memakai
kalibrasi lama diselaraskan ke \HopperFull{} dan \HopperEmpty{} dengan hanya
menyentuh baris yang masih persis bernilai bawaan lama; dan tabel \code{commands}
dibangun ulang bila batasan tipenya belum mengenal perintah uji servo. Ketiganya
menjaga basis data lama tetap terpakai tanpa langkah manual.
(\file{dashboard/db.py})

\section{Kolom yang Ada tetapi Tidak Dipakai}

Tiga kolom layak disebut karena keberadaannya dapat menyesatkan pembaca skema.
\begin{itemize}[leftmargin=*,itemsep=2pt]
  \item \code{feeding_logs.schedule_id} tidak pernah diisi; jalur penulisan riwayat
        terjadwal menuliskan null, sehingga nama jadwal tidak dapat ditampilkan
        lewat join untuk riwayat terjadwal.
  \item \code{feeding_logs.distance_before_cm} tidak pernah diisi; jalur penulisan
        mengisi persentase sebelum, bukan jarak sebelum.
  \item \code{commands.status} mengizinkan nilai \code{pending} dan sapuan timeout
        memeriksanya, tetapi perintah selalu ditulis dengan status terkirim,
        sehingga \code{pending} tidak pernah muncul pada data nyata.
\end{itemize}

Ketiganya bukan kesalahan yang menghentikan sistem; semuanya dicatat sebagai temuan
pada \chapref{ch:09-catatan-asumsi} agar tidak dianggap perilaku yang disengaja.
